\documentclass[12pt]{article}

\usepackage{amsmath}
\usepackage{german}
\usepackage{graphicx}
\DeclareMathOperator\sgn{sgn}

%\graphicspath{{/home/ritz/Teaching/fcn/fig/}}

%----------------- layout ----------------------------------------

\oddsidemargin=0in
\textwidth=6in
\topmargin=-0.7in
\textheight=9.5in
\baselineskip=16pt
\font\mittel=cmr17
\font\roman=cmr12
\font\klein=cmr9
\font\mgross=cmr10
\newfont{\CAP}{cmcsc10 scaled 900}
\pagestyle{empty}
\parindent=0mm

%----------------- macros ----------------------------------------
\newcommand{\conc}[1]{[{\rm C}]_{\text{#1}}}
\newcommand{\units}[1]{\text{#1}}
\newcommand{\K}{\text{K}}
\newcommand{\Na}{\text{Na}}
\newcommand{\Cl}{\text{Cl}}
\newcommand{\Leak}{\text{leak}}
\newcommand{\water}{\text{H}_2\text{O}}
\newcommand{\tderiv}[1]{\frac{\mathrm{d}#1}{\mathrm{d}t}}
%-----------------------------------------------------------------


\begin{document}

\parbox{2cm}{
\includegraphics[width=1.8cm]{../bccnlogo.pdf}
}
\parbox{11cm}{
\begin{center}
\large HUMBOLDT-UNIVERSIT"AT \hskip 0.1 cm ZU \hskip 0.1 cm BERLIN
\vskip 0.1 true cm
\mgross BERNSTEIN CENTER FOR\\ COMPUTATIONAL NEUROSCIENCE
\end{center}
}
\parbox{2cm}
{
\hfill
\includegraphics[width=1.8cm]{../hublogo.pdf}
}

\vskip 0.6 true cm

\leftline{\CAP Humboldt-Universit"at zu Berlin 
\hfill Phone: 030/2093-9110}
\leftline{\CAP Philippstr. 13 House 6
             \hfill Fax: 030/2093-6771}
\leftline{\CAP \hfill webpage: http://www.bccn-berlin.de/}


\vskip 0.8 true cm
\centerline{\bf Models of Neural Systems, WS 2010/11}
\centerline{\bf Project 2: Spike-timing-dependent plasticity}
\centerline{Project presentation and report submission: February,
15th, 2011}

\vskip 0.8 true cm

{\bf Background}

Learning in the brain is believed to be implemented by synapses that are
able to adjust their weights to different input and output
distributions. However, the mechanisms underlying this synaptic
plasticity still remain unclear. Hebb postulated that the increase in
synaptic weights requires coherent pre- and post-synaptic activity,
but he did not provide much details on how these activities could be
defined. Here, you will study a plasticity rule, that is based on the relative
timing of spikes in pre- and post-synaptic terminals. You will test if such
a rule will lead to a stable distribution of synaptic weights and the
generation of output spike times with realistic statistics.

\vskip 0.8 true cm

{\bf Problems}

\begin{enumerate}
    \item \textit{Literature review.} Study the paper of Song, Miller,
        Abbott, Nature Neuroscience 2000. Focus on the figures and the method
        section. What is the hypothesis? What are the main results
        and conclusions? What are the main assumptions of the model?
        What are the strengths and weaknesses of the paper?

    \item
        Implement the leaky integrate-and-fire model with inhibitory
        and excitatory inputs described in
        the paper (see Methods section):

        \begin{equation}
            \tau_m
            \tderiv{V}=V_\mathrm{rest}-V+g_\mathrm{exc}(t)(E_{exc}-V)+g_\mathrm{inh}(t)(E_\mathrm{inh}-V)
            \label{eq:lif}
        \end{equation}
    \item Assume that each presynaptic spike at synapse $a$ results in a sudden
        change of the post-synaptic conductance, which then decays
        exponentially until the next input arrives:
        \begin{equation}
            g_\mathrm{syn} \rightarrow g_\mathrm{syn} +
            \bar{g}_\mathrm{a},\;\text{if } t =
            t_\mathrm{pre}
            \label{eq:conductance2}
        \end{equation}
        \begin{equation}
            \tau_\mathrm{syn} \tderiv{g_\mathrm{syn}} =
            -g_\mathrm{syn},
            \label{eq:conductance2}
        \end{equation}
        where $\mathrm{syn} \in  \left\{\mathrm{exc},
        \mathrm{inh}\right\}$, $t_\mathrm{pre}$ is the time of a pre-synaptic
        time and $\bar{g}_\mathrm{a}$ is the peak synaptic conductance of
        synapse $a$.
    \item  Spike-timing-dependent synaptic plasticity is implemented by
        means of modification of the peak conductances
        of excitatory synapses $\bar{g}_\mathrm{a}$ (inhibitory
        synapses are not changed) according
        to the rules:
        \begin{equation} 
            \bar{g}_\mathrm{a} \rightarrow
            \bar{g}_\mathrm{a}+M(t)\bar{g}_\mathrm{max},\; \;
            M \rightarrow M - A_-, \; \;
            \text{if } t =
            t_\mathrm{pre}
        \end{equation}
        \begin{equation} 
            \bar{g}_\mathrm{a} \rightarrow
            \bar{g}_\mathrm{a}+
            P_a(t)\bar{g}_\mathrm{max},\; \;
            P_a \rightarrow P_a + A_+, \; \;
            \text{if } t =
            t_\mathrm{post}
        \end{equation}
        %Every time the post-synaptic neuron fires an action potential,
        %$M(t)$ is decremented by an amount $A_-$, and every time a
        %synapse receives an action potential from the pre-synaptic
        %neuron, $P(t)$ is incremented by an amount $A_{+}$.
        If $\bar{g}_\mathrm{a}<0$ the peak conductance is clamped to zero and
        if $\bar{g}_\mathrm{a}>\bar{g}_\mathrm{max}$ it is reset to
        $\overline{g}_\mathrm{max}$. If there are no pre- or
        post-synaptic spikes $M(t)$ and $P(t)$ decay exponentially:
        \begin{equation}
            \tau_{-}\tderiv{M}=-M \text{ and } \tau_{+}\tderiv{P_a}=-P_a.
            \label{eq:stdp}
        \end{equation}


    \item For the initial excitatory peak conductances take
        $\bar{g}_\mathrm{a}=\bar{g}_\mathrm{max}$ for all $a$. Simulate the
        model with $N=1000$ excitatory and $M=200$ inhibitory synapses,
        both of which receive independent
        Poisson spike trains (10 Hz for inhibitory, 10 -- 40 Hz
        for excitatory). Plot the final weight distribution for
        different excitatory firing rates. Plot the output firing rate
        and coefficient of variation as a function of the input firing
        rate.

    \item Simulate the model with inputs consisting of bursts of
        action potentials. To this end, take the inputs to be silent
        except for isolated events represented by bursts of spikes
        with a Poisson distribution at 100 Hz for 20 ms. The inputs
        arrive at each synapse with random latencies drawn from a
        Gaussian distribution (mean 0, standard deviation 15 ms). Plot
        a response to the same burst of spikes at the beginning and
        at the end of the simulation. Explain the observed differences.
        Plot the weights of synapses as a function of the relative
        latencies of incoming spike bursts (Song et al., 2000, Figure 4).
\end{enumerate}

\vskip 0.8 true cm
{\bf Literature}
\vskip 0.8 true cm

S.~Song, K.D.~Miller and L.F.~Abbott (2000). Competitive Hebbian learning
through spike-timing-dependent synaptic plasticity.
\emph{Nature Neuroscience} 3(9):919-926.\newline

\vfill
\centerline{\CAP Contact}
\CAP

\begin{tabular}{lll}
\underline{Urs Bergmann} & Phone: 2093-8924 & Email:
urs.bergmann(at)biologie.hu-berlin.de \\
Matthias Guggenmos & & Email: matthias.guggenmos(at)bccn-berlin.de \\
Richard Kempter \hfill & Phone: 2093-8925 \hfill & Email:
r.kempter(at)biologie.hu-berlin.de \\
\end{tabular}

\end{document}
